\medskip

On considère la suite $\left(u_{n}\right)$ définie par: $\left\{\begin{array}{l c l}
u_{1}&=&\dfrac{1}{\e} \\
u_{n+1}&=&\dfrac{1}{\e}\left(1+\dfrac{1}{n}\right) u_{n}\: \text {pour tout entier } n \geqslant 1.\end{array}\right.$

\begin{enumerate}
	\item Calculer les valeurs exactes de $u_{2}$ et $u_{3}$. On détaillera les calculs.

	\item On considère une fonction écrite en langage Python qui, pour un entier naturel $n$ donné, affiche le terme $u_{n}$. Compléter les lignes $L_{2}$ et $L_{4}$ de ce programme.

	\begin{center}
		\begin{tabular}{|c|>{\ttfamily}l|}
			\hline
			$L_{1}$ & def suite(n): \\
			$L_{2}$ & ~~\ldots \ldots \ldots \ldots \ldots \ldots \\
			$L_{3}$ & ~~for i in range(1, n): \\
			$L_{4}$ & ~~~~u=\ldots \ldots \ldots \ldots \ldots \ldots\\
			$L_{5}$ & return u \\
			\hline
		\end{tabular}
	\end{center}

	\item On admet que tous les termes de la suite $\left(u_{n}\right)$ sont strictement positifs.
	\begin{enumerate}
		\item Montrer que pour tout entier naturel $n$ non nul, on a : $1+\dfrac{1}{n} \leqslant \e$.

		\item En déduire que la suite $\left(u_{n}\right)$ est décroissante.
		\item La suite $\left(u_{n}\right)$ est-elle convergente ? Justifier votre réponse.
	\end{enumerate}

	\item \begin{enumerate}
		\item Montrer par récurrence que pour tout entier naturel non nul, on a : \quad $u_{n}=\dfrac{n}{\e^{n}}$.
		\item En déduire, si elle existe, la limite de la suite $\left(u_{n}\right)$.
	\end{enumerate}
\end{enumerate}
